\section{Intelligence per Watt：把能耗纳入能力扩展}

\subsection{从云端 mainframe 到 local AI}

\videofigure{figures/intelligence-per-watt-title.jpg}{课程后半提出 intelligence per watt，用单位功耗可提供的模型能力衡量本地 AI 效率。}{00:40:08--00:42:00}

\videofigure{figures/mainframe-era.jpg}{当前 AI 类似 mainframe 时代：大部分推理由集中云端基础设施提供，而需求增长快于供给。}{00:43:48--00:45:22}

课程观察到大量真实 chatbot query 只需要写作、信息查询、一般建议和常见 reasoning，并不总需要 frontier cloud model。若小于约 20B 参数的本地模型能处理多数请求，就可以降低网络延迟、云端成本与隐私暴露。

\videofigure{figures/local-model-progress.jpg}{本地小模型的能力随模型和设备代际快速增长，为计算从云端部分迁移到端侧创造条件。}{00:45:22--00:46:12}

\videofigure{figures/local-inference-question.jpg}{核心系统问题是：哪些推理应本地执行，哪些请求值得升级到云端 frontier model。}{00:46:12--00:47:18}

\begin{importantbox}{Local/cloud 不是二选一，而是动态路由}
简单、隐私敏感、低延迟请求优先本地；需要大模型知识、长 context 或昂贵工具的请求升级云端。路由器本身要估计难度、风险和资源。
\end{importantbox}

\subsection{IPW 的定义}

\videofigure{figures/intelligence-per-watt-metric.jpg}{Intelligence per watt 将可解决请求比例与平均功耗结合，比较不同模型--硬件组合。}{00:47:18--00:48:18}

课程使用的直觉指标可写为：

$$
\operatorname{IPW}(m,h)=
\frac{C(m)}{\overline P(m,h)},
$$

\begin{itemize}
    \item $m$：模型；
    \item $h$：运行硬件；
    \item $C(m)$：模型对目标 query workload 的可解决比例或能力分数；
    \item $\overline P(m,h)$：该模型在硬件上的平均功耗；
    \item IPW 越高，单位功耗提供的可用智能越多。
\end{itemize}

\videofigure{figures/efficiency-study.jpg}{研究跨多个本地模型、设备与真实 query workload 测量能力、延迟、准确率和能耗。}{00:48:18--00:49:04}

\begin{warningbox}{IPW 必须绑定 workload 和质量阈值}
若一个模型只擅长简单问答，它可能功耗极低却无法处理复杂请求。比较时必须固定查询分布、输出质量标准、context 长度和 batch 设置。
\end{warningbox}

\subsection{两年 5.3 倍效率提升}

\videofigure{figures/local-ai-findings.jpg}{本地模型可准确处理约 88.7\% 的研究 query；两年内模型能力提高约 3.1 倍，硬件效率提高约 1.7 倍，IPW 合计提高约 5.3 倍。}{00:49:04--00:52:18}

课程同时指出 Apple M4 Max 的 IPW 仍约比 B200 低 1.5 倍，因为企业 GPU 专门针对 LLM workload 优化，而本地芯片还要服务图形、媒体和通用应用。未来端侧 accelerator、kernel 和模型架构仍有巨大优化空间。

\begin{knowledgebox}{5.3 倍近似来自软硬件乘法}
$3.1\times1.7\approx5.27$。模型压缩、训练方法与架构提高 numerator，编译器、kernel、内存系统和芯片提高 denominator；联合优化通常比只优化一端更有效。
\end{knowledgebox}

\subsection{本章小结}

Intelligence per watt 把“更强模型”改写为“在真实 workload 上用更少能量提供足够能力”。Local AI 的快速进步支持混合推理，但硬件专用化和任务路由仍是关键。

